We study the problem of efficient verifier allocation during test-time compute scaling
for large language models. Recent work has established that process reward models (PRMs)
used as verifiers are asymptotically optimal for scaling test-time compute; however,
invoking a verifier at every node of a search tree incurs costs that are linear in the
search budget, limiting scalability. We introduce \emph{Dynamic Verifier Allocation
Monte Carlo Tree Search} (DVA-MCTS), an algorithm that adaptively decides when to
invoke the verifier based on a budget-sensitive uncertainty threshold. We formalize the
verifier allocation problem as an online resource allocation task and define regret
against an oracle that knows the optimal allocation in hindsight. Under a
\emph{Lipschitz score continuity} assumption on the PRM---which we validate
empirically---we prove that DVA-MCTS achieves cumulative regret
$R(T) = O\!\left(\sqrt{T \log K}\right)$ against the oracle while invoking the verifier
at most $N_{\mathcal{T}} = K(K^D-1)/(K-1)$ times in total---a constant independent of~$T$---so
the call fraction $\mathcal{B}/T \to 0$ as $T \to \infty$.
We further prove a matching $\Omega(\sqrt{T})$ lower bound, establishing rate-optimality.
Empirically, DVA efficiency unfolds across three regimes: an \emph{exploration}
regime ($B \ll N_{\mathcal{T}}$) where all algorithms are statistically equivalent
($p{>}0.10$, Wilcoxon); a \emph{transition} regime ($B \lesssim N_{\mathcal{T}}$)
where the Lipschitz proxy mechanism delivers $5$--$73\%$ savings; and an
\emph{exploitation} regime ($B \gg N_{\mathcal{T}}$) where the single-verification
ceiling yields savings exceeding $90\%$.
The exploitation regime is demonstrated at two scales:
$\mathbf{93.7\%}$ savings at $B{=}16{,}384$ ($K{=}2$, $D{=}12$, $N_{\mathcal{T}}{=}8{,}190$);
and $\mathbf{94.6\%}$ savings at $B{=}400$ for the directly-checkable small-tree
setting ($K{=}2$, $D{=}4$, $N_{\mathcal{T}}{=}30$), where DVA call counts plateau at
$21.5 < N_{\mathcal{T}}$, directly confirming the call-ceiling theorem.
An adaptive $\hat{L}$ estimator---tracking the running maximum of observed
score-change ratios---requires no advance knowledge of $L$ and achieves the lowest
empirical regret.
On $100$ real GSM8K problems (Math-Shepherd PRM) with binary labels, DVA-MCTS
reduces calls by $8.8\%$ while improving accuracy by $3.6$ pp.
Converting binary labels to running-average continuous scores
($L_{\text{eff}}{=}0.155$) recovers $\mathbf{49.9\%}$ call savings, confirming
that $L_{\text{eff}}$ is the primary driver: neural regression-trained PRMs with
$L_{\text{eff}} \lesssim 0.1$ are predicted to exceed $73\%$ savings.
